\chapter{Requisitos e metodologia do projeto}

\section{Requisitos funcionais}

Para a proposta da BMS 4S, consideram-se como requisitos funcionais:

\begin{itemize}
  \item leitura da tensão individual das quatro células;
  \item leitura de corrente do pack;
  \item leitura de temperatura por sensores NTC;
  \item controle independente dos caminhos de carga e descarga;
  \item suporte a balanceamento passivo;
  \item disponibilização de telemetria do pack.
\end{itemize}

\section{Requisitos de segurança}

Como requisitos de segurança, destacam-se:

\begin{itemize}
  \item bloqueio por sobretensão;
  \item bloqueio por subtensão;
  \item bloqueio por sobrecorrente;
  \item resposta segura em condição de curto-circuito;
  \item bloqueio por sobretemperatura;
  \item comportamento \textit{fail-safe} em caso de falha crítica.
\end{itemize}

\section{Parâmetros de configuração}

Os parâmetros iniciais do sistema são organizados em uma estrutura de configuração contendo:

\begin{itemize}
  \item características do pack;
  \item parâmetros de medição;
  \item limites de proteção;
  \item estratégia de balanceamento;
  \item convenções de telemetria.
\end{itemize}

\section{Metodologia de desenvolvimento}

A metodologia adotada foi composta por:

\begin{enumerate}
  \item levantamento bibliográfico;
  \item análise de AFEs e arquiteturas de mercado;
  \item definição dos requisitos do sistema;
  \item modelagem da arquitetura de hardware;
  \item modelagem da máquina de estados;
  \item implementação da base inicial de firmware.
\end{enumerate}
